\documentclass[a4paper,11pt]{article}
\usepackage[utf8]{inputenc}
\usepackage[T1]{fontenc}
\usepackage[brazil]{babel}
% Ajuste de geometria para acomodar o cabeçalho grande
\usepackage[left=2cm, right=2cm, top=3.5cm, bottom=2.5cm, headheight=50pt]{geometry}
\usepackage{xcolor}
\usepackage{graphicx}
\usepackage{tcolorbox}
\usepackage{helvet}
\usepackage{tikz}
\usepackage{fancyhdr}
\usepackage{array}
\usepackage{colortbl}
\usepackage{hyperref}

% --- CONFIGURAÇÕES DE LAYOUT E BIBLIOTECAS ---
\tcbuselibrary{skins, raster, breakable}
\renewcommand{\familydefault}{\sfdefault}
\setlength{\parindent}{0pt}
\setlength{\parskip}{0.8em}

% --- PALETA DE CORES KROOA (ALTO CONTRASTE) ---
\definecolor{krooaNavy}{RGB}{12, 45, 60}      % Fundo Escuro
\definecolor{krooaLime}{RGB}{205, 239, 94}    % Acento Neon (Verde Lima)
\definecolor{krooaRed}{RGB}{231, 76, 60}      % Vermelho Alerta
\definecolor{krooaOrange}{RGB}{243, 156, 18}  % Laranja (Atenção/Suave)
\definecolor{krooaGray}{RGB}{240, 242, 245}   % Fundo Suave das Tabelas
\definecolor{textDark}{RGB}{44, 62, 80}       % Texto Leitura

% --- CABEÇALHO E RODAPÉ PERSONALIZADOS ---
\pagestyle{fancy}
\fancyhf{} % Limpa tudo
\renewcommand{\headrulewidth}{0pt}

% Desenho do Cabeçalho (Barra Azul no Topo)
\fancyhead[C]{%
    \begin{tikzpicture}[remember picture, overlay]
        \fill[krooaNavy] (current page.north west) rectangle ([yshift=-3cm]current page.north east);
        \node[anchor=west, text=white, font=\bfseries\Huge] at ([xshift=2cm, yshift=-1.5cm]current page.north west) {KROOA \textcolor{krooaLime}{INTELLIGENCE}};
        \node[anchor=east, text=white, align=right, font=\small] at ([xshift=-2cm, yshift=-1.5cm]current page.north east) {
            \textbf{RELATÓRIO DE AUDITORIA FORENSE}\\
            PROTOCOLO: SHOPPER-GO-01 \\
            DATA: \today
        };
        % Linha de destaque verde
        \fill[krooaLime] (current page.north west) ++(0,-2.9cm) rectangle ++(\paperwidth, 0.1cm);
    \end{tikzpicture}
}

% Desenho do Rodapé
\fancyfoot[C]{%
    \begin{tikzpicture}[remember picture, overlay]
        \fill[krooaGray] (current page.south west) rectangle ([yshift=1.5cm]current page.south east);
        \node[text=krooaNavy, font=\footnotesize] at ([yshift=0.75cm]current page.south) {Confidencial - Gerado por Krooa AI Engine - Página \thepage};
    \end{tikzpicture}
}

% --- ÍCONES MANUAIS ---
\newcommand{\iconAlert}{%
\begin{tikzpicture}[baseline=-3pt]
    \fill[krooaOrange] (0,0) circle (0.25);
    \node[white, font=\bfseries\small] at (0,0) {!};
\end{tikzpicture}%
}

\newcommand{\iconCheck}{%
\begin{tikzpicture}[baseline=-3pt]
    \fill[krooaLime] (0,0) circle (0.25);
    \draw[krooaNavy, very thick] (-0.1,-0.05) -- (0,0.1) -- (0.15,-0.15);
\end{tikzpicture}%
}

\begin{document}

% --- BADGE DE STATUS ---
\begin{flushright}
    \begin{tikzpicture}
        \node[fill=krooaOrange!10, draw=krooaOrange, rounded corners, inner sep=5pt, line width=1pt] {
            \textbf{\textcolor{krooaOrange}{STATUS: OPORTUNIDADE DE MELHORIA DETECTADA}}
        };
    \end{tikzpicture}
\end{flushright}

% --- SEÇÃO 1: O CENÁRIO ---
\begin{tcolorbox}[
    enhanced,
    title=\textbf{1. O CENÁRIO DO TESTE (A ISCA)},
    colframe=krooaNavy, colback=white, coltitle=white,
    boxrule=0.5pt, drop shadow
]
    \textbf{Script Aplicado:} "O Viajante com Pressa" \\
    \textbf{Perfil do Avatar:} Executivo(a) de Goiás (DDD 62) em viagem rápida de negócios ou visita familiar. \\
    \textbf{Cenário:} O "High Ticket" de Emergência Estética (Horário de Pico) \\
    \textbf{Contexto:} Quebrou uma faceta/dente ou caiu um provisório - situação estética que gera desespero. \\
    \textbf{Gatilhos de Urgência:} "Lascou pedacinho do dente da frente" + "Reuniões importantes" + "Fico só até depois de amanhã". \\
    \textbf{Potencial Financeiro:} "Pago particular se precisar" + Disponibilidade para resolver urgente.
\end{tcolorbox}

% --- SEÇÃO 2: A LINHA DO TEMPO ---
\vspace{0.5cm}
\section*{\textcolor{krooaNavy}{2. ANÁLISE MINUTO A MINUTO}}

% Tabela customizada com tcolorbox
\begin{tcolorbox}[enhanced, frame hidden, colback=white]
    \renewcommand{\arraystretch}{1.5}
    \begin{tabular}{ | p{0.12\textwidth} | p{0.20\textwidth} | p{0.58\textwidth} | }
        \hline
        \rowcolor{krooaNavy} 
        \textbf{\textcolor{white}{HORA}} & \textbf{\textcolor{white}{AGENTE}} & \textbf{\textcolor{white}{ANÁLISE DE FLUXO}} \\
        \hline
        
        \textbf{17:49} & Paciente (Lead) & 
        \small Disparo inicial relatando dor e urgência de tempo. \\
        \hline
        
        \textbf{17:49} & Clínica (Humano) & 
        \textbf{Tempo de Resposta: < 1 Min.} (Excelente velocidade inicial). \newline
        \textit{"Kibom!!! Temos sim!!! qual seu nome completo e cpf?"} \\
        \hline
        
        \multicolumn{3}{c}{
            \begin{tcolorbox}[colback=krooaOrange!5, colframe=krooaOrange, boxrule=0.5pt, arc=0mm, left=2mm, right=2mm]
                \textbf{\textcolor{krooaOrange}{PONTO DE OTIMIZAÇÃO \#1: ALINHAMENTO DE TOM}} \\
                \footnotesize
                1. \textbf{Comunicação:} A equipe demonstrou muita proatividade. Porém, o uso de "Kibom!!!" para uma dor física pode soar informal. Sugerimos um tom mais consultivo.\\
                2. \textbf{Momento do Cadastro:} Pedir o CPF antes de acolher a dor cria uma fricção desnecessária. O ideal é confirmar o horário primeiro.
            \end{tcolorbox}
        } \\
        
        \textbf{17:50} & Paciente & 
        \small Nega passar dados e insiste na confirmação. \\
        \hline
        
        \textbf{17:51} & Clínica (Humano) & 
        \textit{"kl pena, tenho um encaixe amanhã..."} \\
        \hline
        
        \multicolumn{3}{c}{
            \begin{tcolorbox}[colback=krooaOrange!5, colframe=krooaOrange, boxrule=0.5pt, arc=0mm, left=2mm, right=2mm]
                \textbf{\textcolor{krooaOrange}{PONTO DE OTIMIZAÇÃO \#2: PRECISÃO NA DIGITAÇÃO}} \\
                \footnotesize
                O erro \textbf{"kl pena"} é um sintoma clássico de \textit{multitarefa}. Sua equipe provavelmente estava atendendo alguém presencialmente enquanto digitava. Isso comprova sobrecarga operacional, não má vontade.
            \end{tcolorbox}
        } \\
    \end{tabular}
\end{tcolorbox}

% --- SEÇÃO 3: DIAGNÓSTICO VISUAL ---
\vspace{0.2cm}
\begin{tcolorbox}[
    enhanced, 
    colback=krooaGray, 
    colframe=white, 
    title=\textcolor{krooaNavy}{\textbf{3. DIAGNÓSTICO OPERACIONAL}},
    coltitle=krooaNavy,
    fonttitle=\large\bfseries,
    frame hidden
]
    \begin{itemize}
        \item[\iconAlert] \textbf{Dependência Humana:} A resposta rápida (17:49) mostra que sua equipe é atenta, mas depende da disponibilidade momentânea. Se estivessem ao telefone, o lead esperaria.
        \item[\iconAlert] \textbf{Refinamento Premium:} Para tratamentos de alto valor, a precisão da escrita e o acolhimento inicial são vitais para transmitir a mesma qualidade técnica do Doutor.
    \end{itemize}
\end{tcolorbox}

% --- SEÇÃO 4: A SOLUÇÃO KROOA (ATUALIZADA) ---
\vspace{0.5cm}

\begin{tcolorbox}[
    skin=enhanced,
    colback=krooaNavy,
    colframe=krooaLime,
    boxrule=1pt,
    arc=4mm,
    shadow={2mm}{-2mm}{0mm}{black!30},
    title=\textbf{\large KROOA: INTELIGÊNCIA CLÍNICA VERTICAL},
    coltitle=krooaLime,
    fonttitle=\bfseries
]
    \color{white}
    Diferente de "Chatbots" que apenas respondem, o Krooa é um \textbf{Sistema Vertical} que atua em todas as camadas do processo clínico:

    \vspace{0.3cm}

    \begin{itemize}
        \item[\iconCheck] \textbf{Camada de Comunicação:} Elimina erros humanos ("kl pena"), garantindo um padrão de escrita "Concierge" e negociação persuasiva.
        \item[\iconCheck] \textbf{Camada de Processo:} Inverte a lógica da burocracia. O sistema fecha a agenda emocionalmente antes de pedir dados cadastrais, aumentando a conversão.
        \item[\iconCheck] \textbf{Camada de Infraestrutura:} Atua 24/7 de forma sistêmica, escalando o atendimento sem depender da disponibilidade física da recepção.
    \end{itemize}

    \vspace{0.5cm}
    \centering
    \begin{tcolorbox}[colback=krooaLime, colframe=krooaLime, width=12cm, halign=center, arc=4mm, boxsep=15pt, top=8pt, bottom=8pt]
        \href{https://krooa.com}{\textbf{\Large\textcolor{krooaNavy}{PROFISSIONALIZAR MEU SISTEMA}}}
    \end{tcolorbox}
    \vspace{0.3cm}
\end{tcolorbox}

\end{document}